% =============================================================================
\chapter{Profile Export Pipeline}
\label{ch:export}
% =============================================================================

After optimisation, the variable-thickness element mesh must be converted
into a clean CAD-ready profile.  The \texttt{exportProfile} function uses an
image-processing approach that is topology-agnostic: it handles beams,
trusses, branching structures, and any geometry with holes or disconnected
regions.

This chapter describes the algorithm in detail.

% -----------------------------------------------------------------------------
\section{Algorithm Overview}
\label{sec:export:overview}
% -----------------------------------------------------------------------------

The pipeline has six stages:

\begin{enumerate}
    \item \textbf{Filter.}  Optionally remove elements thinner than a
          fraction of the maximum thickness.
    \item \textbf{Rasterise.}  Draw every active element as a filled
          rectangle onto a high-resolution binary image.
    \item \textbf{Clean.}  Apply morphological operations to close gaps
          and remove noise.
    \item \textbf{Trace.}  Extract all closed boundary contours using
          \texttt{bwboundaries}.
    \item \textbf{Fit.}  Convert pixel contours to world coordinates and
          fit a smoothing spline to each loop.
    \item \textbf{Export.}  Write each spline loop as a text file of
          $[X,\,Y,\,Z]$ points for SolidWorks import.
\end{enumerate}

% -----------------------------------------------------------------------------
\section{Rasterisation}
\label{sec:export:raster}
% -----------------------------------------------------------------------------

Each active element is represented by a rectangle centred on the element
axis with width equal to the element thickness.  For an element connecting
nodes $(x_1, y_1)$ and $(x_2, y_2)$ with thickness~$t$ and orientation
$\alpha = \text{atan2}(y_2 - y_1,\, x_2 - x_1)$, the four corners are:

\begin{equation}
    \begin{bmatrix} x_k \pm p_x \\ y_k \pm p_y \end{bmatrix}, \qquad
    p_x = -\tfrac{t}{2}\sin\alpha, \quad
    p_y = \phantom{-}\tfrac{t}{2}\cos\alpha, \quad
    k \in \{1, 2\}.
\end{equation}

These corners are mapped to pixel coordinates and rasterised using MATLAB's
\texttt{poly2mask}, which fills the interior of each polygon.  All element
masks are combined with a logical OR.

The image resolution is controlled by the \texttt{resolution} parameter
(pixels per metre, default 50\,000).  To prevent excessive memory use, the
image dimensions are capped at $8000\times8000$ pixels.

% -----------------------------------------------------------------------------
\section{Morphological Cleanup}
\label{sec:export:morph}
% -----------------------------------------------------------------------------

Adjacent elements may leave one- or two-pixel gaps at joints due to
discretisation.  A morphological closing operation (\texttt{imclose}) with a
small disk structuring element fills these gaps:

\begin{equation}
    B_{\text{clean}} = (B \oplus S) \ominus S,
\end{equation}

where $S$ is a disk with radius proportional to $10\%$ of the maximum
element thickness.  This bridges small gaps without significantly altering
the outline shape.

Very small isolated pixel clusters (noise) are removed with
\texttt{bwareaopen}.

% -----------------------------------------------------------------------------
\section{Boundary Extraction}
\label{sec:export:boundary}
% -----------------------------------------------------------------------------

MATLAB's \texttt{bwboundaries} traces all closed contours in the binary
image using an 8-connected Moore neighbourhood.  It returns:

\begin{itemize}[nosep]
    \item \textbf{Outer boundaries} --- one per connected region.
    \item \textbf{Hole boundaries} --- one per enclosed void (e.g.\ the
          interior of a truss after thin elements are filtered).
\end{itemize}

Each boundary is returned as a list of $[\text{row},\, \text{col}]$ pixel
indices.  These are converted back to world coordinates using the known
image bounds:

\begin{align}
    x_{\text{world}} &= \frac{c - 1}{W - 1}\,(x_{\max} - x_{\min}) + x_{\min}, \\[3pt]
    y_{\text{world}} &= \frac{H - r}{H - 1}\,(y_{\max} - y_{\min}) + y_{\min},
\end{align}

where $(r, c)$ are the row and column pixel indices, and $W$, $H$ are the
image width and height.

% -----------------------------------------------------------------------------
\section{Spline Fitting}
\label{sec:export:spline}
% -----------------------------------------------------------------------------

Each boundary loop is parameterised by cumulative arc length~$s$.  After
removing duplicate consecutive points, a smoothing cubic spline is fitted
using \texttt{csaps} (Curve Fitting Toolbox) or, if unavailable,
\texttt{pchip} interpolation:

\begin{equation}
    x(s) = \text{csaps}(s_i,\, x_i,\, p,\, s_{\text{fine}}), \qquad
    y(s) = \text{csaps}(s_i,\, y_i,\, p,\, s_{\text{fine}}),
\end{equation}

where $p$ is the smoothing parameter (default $0.999999999 \approx 1$,
giving near-interpolation with minimal smoothing) and $s_{\text{fine}}$ is
a uniform grid with \texttt{num\_spline\_points} samples.

The spline is closed by appending the first point to the end.

% -----------------------------------------------------------------------------
\section{File Format}
\label{sec:export:format}
% -----------------------------------------------------------------------------

Each loop is exported as a plain-text file with one point per line in the
format:

\begin{verbatim}
    X.XXXXXX Y.YYYYYY 0.000000
\end{verbatim}

The $Z$~coordinate is zero (2-D planar structure).  Coordinates are in the
units specified by \texttt{scale\_to\_mm} (default: millimetres).

Files are named \texttt{<base>\_loop\_<n>.txt} for outer boundaries and
\texttt{<base>\_hole\_<n>.txt} for interior voids.  A \texttt{.mat} file
containing the full \texttt{profile} struct is also saved.

% -----------------------------------------------------------------------------
\section{SolidWorks Import Workflow}
\label{sec:export:solidworks}
% -----------------------------------------------------------------------------

To import the exported profiles into SolidWorks:

\begin{enumerate}
    \item Open a new part and select the Front Plane.
    \item \textbf{Insert} $\to$ \textbf{Curve} $\to$
          \textbf{Curve Through XYZ Points}.
    \item Browse to the exported \texttt{.txt} file and click \textbf{OK}.
    \item Repeat for each loop and hole file.
    \item Create a new sketch on the Front Plane and
          \textbf{Convert Entities} to pull the curves into the sketch.
    \item Close any open loop endpoints with short sketch lines if needed.
    \item \textbf{Extrude Boss/Base} the outer boundary to the desired width.
    \item \textbf{Extrude Cut} any hole boundaries through the part.
\end{enumerate}

\begin{notebox}[Tip]
For beam-type profiles with a single outer loop and no holes, a single
\textbf{Extrude Boss/Base} produces the finished part in one operation.
For truss-type profiles, import the outer boundary first, extrude, then
cut the holes.
\end{notebox}

% -----------------------------------------------------------------------------
\section{Controlling Quality}
\label{sec:export:quality}
% -----------------------------------------------------------------------------

The key parameters affecting export quality are:

\begin{table}[H]
\centering
\caption{Profile export quality parameters.}
\begin{tabularx}{\textwidth}{lX}
\toprule
\textbf{Parameter} & \textbf{Effect} \\
\midrule
\texttt{resolution}           & Higher values capture thinner elements and
                                sharper corners, at the cost of memory and
                                processing time. \\
\texttt{smoothing}            & Values very close to~1 (e.g.\ $0.999999999$)
                                preserve fine detail; lower values
                                (e.g.\ $0.95$) produce smoother, rounder
                                profiles. \\
\texttt{min\_thickness\_frac} & For trusses, setting this to $0.10$--$0.20$
                                removes vanished material and leaves clean
                                load-path outlines. \\
\texttt{num\_spline\_points}  & More points give a denser export file,
                                which may improve CAD import fidelity for
                                complex shapes. \\
\bottomrule
\end{tabularx}
\end{table}